% Part:sets-functions-relations
% Chapter: size-of-sets
% Section: comparing-sizes

\documentclass[../../../include/open-logic-section]{subfiles}

\begin{document}

\olfileid{sfr}{siz}{car}

\olsection{不同大小的集合与Cantor定理}

\begin{explain}
我们已经给出了两个集合“大小相同”这一概念的精确定义。
我们同样可以给出“一个集合比另一个小”这一概念的精确定义。
我们对“小于（或等势）”的定义，需要的不是两个集合之间的双射，
而是从第一个集合到第二个集合的单射。
如果这样的函数存在，那么第一个集合的大小就小于或等于第二个集合的大小。
直观上，从一个集合到另一个集合的单射保证了该函数的值域至少和定义域包含同样多的元素，
因为定义域中没有两个元素映射到值域中的同一个元素。
\end{explain}

\begin{defn}
$A$ \emph{不大于}~$B$，记作 $\cardle{A}{B}$，当且仅当存在单射 $f \colon A \to B$。
\end{defn}

显然，这是一个自反且传递的关系，但不是对称的（留作练习）。
我们还可以引入一个概念，用以表明一个集合（严格地）小于另一个集合。

\begin{defn}
$A$ \emph{小于}~$B$，记作 $\cardless{A}{B}$，当且仅当存在单射~$f\colon A \to B$ 但不存在双射~$g\colon A
\to B$，即 $\cardle{A}{B}$ 且 $\cardneq{A}{B}$。
\end{defn}

显然，该关系是非自反且传递的。（这也留作练习。）使用这一记号，我们可以说 $A$ 是可数的当且仅当 $\cardle{A}{\Nat}$，而 $A$ 是不可数的当且仅当 $\cardless{\Nat}{A}$。这使我们可以将
\oliflabeldef{sfr:siz:nen-alt:thm:nonenum-pownat}{%
\olref[sfr][siz][nen-alt]{thm:nonenum-pownat}
重述为如下结论：
$\cardless{\Nat}{\Pow{\Nat}}$}{%
\olref[sfr][siz][nen]{thm:nonenum-pownat}
重述为如下结论：$\cardless{\PosInt}{\Pow{\PosInt}}$}。
事实上，\citet{Cantor1892} 证明了最后这一点具有\emph{完全的一般性}：

\begin{thm}[Cantor]\ollabel{thm:cantor}
对任意集合 $A$，$\cardless{A}{\Pow{A}}$。
\end{thm}

\begin{proof}
映射 $f(x) = \{x\}$ 是一个单射 $f \colon A \to \Pow{A}$，因为若 $x \neq y$，则由外延性原理亦有 $\{x\} \neq \{y\}$，从而 $f(x) \neq f(y)$。因此我们有 $\cardle{A}{\Pow{A}}$。

\begin{editorial}
如果存在 \olref[nen]{sec}，我们给出慢速证明；否则给出与 \olref[nen-alt]{sec} 匹配的快速证明。
\end{editorial}

\oliflabeldef{sfr:siz:nen:sec}{% 
  现在我们证明不可能存在满射函数~$g\colon A \to \Pow{A}$，更不用说双射了，从而 $\cardneq{A}{\Pow{A}}$。
  假设 $g\colon A \to \Pow{A}$。
  由于 $g$ 是全函数，每个 $x \in A$ 都映射到某个子集 $g(x) \subseteq A$。
  我们来证明 $g$ 不可能是满射。为此，定义子集~$\overline{A} \subseteq A$，由其定义可知它不可能在~$g$ 的值域中。令
  \[
  \overline{A} = \Setabs{x \in A}{x \notin g(x)}.
  \]
  由于 $g(x)$ 对所有 $x \in A$ 均有定义，$\overline{A}$ 显然是~$A$ 的一个良定义的子集。
  但它不可能在~$g$ 的值域中。任取 $x \in A$，我们证明 $\overline{A} \neq g(x)$。
  若 $x \in g(x)$，则它不满足 $x \notin g(x)$，因而根据~$\overline{A}$ 的定义，有 $x \notin \overline{A}$。
  若 $x \in \overline{A}$，则它必须满足~$\overline{A}$ 的定义性质，即 $x \in A$ 且 $x \notin g(x)$。
  由于 $x$ 是任取的，这表明对每个 $x \in A$，$x \in g(x)$ 当且仅当 $x \notin \overline{A}$，故 $g(x) \neq \overline{A}$。
  换言之，$\overline{A}$ 不可能在~$g$ 的值域中，这与~$g$ 为满射的假设相矛盾。}{还需证明 $\cardneq{A}{\Pow{A}}$。用归谬法，假设 $\cardeq{A}{\Pow{A}}$，即存在某个双射 $g \colon A \to \Pow{A}$。现在考虑：
  \[
    D = \Setabs{x \in A}{x \notin g(x)}
  \]
  注意 $D \subseteq A$，因此 $D \in \Pow{A}$。由于 $g$ 是双射，存在某个 $y \in A$ 使得 $g(y) = D$。但现在我们有：
  \[
    y \in g(y) \text{ 当且仅当 } y \in D \text{ 当且仅当 } y \notin g(y).
  \]
  得出矛盾；因此 $\cardneq{A}{\Pow{A}}$。}{}
\end{proof}

\begin{explain}
\oliflabeldef{sfr:siz:nen:thm:nonenum-pownat}{将 \olref{thm:cantor} 的证明与
  \olref[nen]{thm:nonenum-pownat} 的证明加以比较颇具启发意义。在那里我们证明了，对于 $\PosInt$ 的任何子集列表 $Z_1$, $Z_2$, \dots，都能构造出一个集合~$\overline{Z}$，保证它不在该列表中。之所以能保证这一点，是因为对于每个 $n \in
  \PosInt$，$n \in Z_n$ 当且仅当 $n \notin \overline{Z}$。这样，总存在某个数，它属于 $Z_n$ 与 $\overline{Z}$ 之一，而不属于另一个。我们在这里遵循同样的思想，只不过现在的索引~$n$ 是~$A$ 的元素而非~$\PosInt$ 的元素。集合 $\overline{A}$ 的定义保证了对每个 $x \in A$，它都不同于~$g(x)$，因为 $x \in g(x)$ 当且仅当 $x \notin \overline{A}$。同样地，总存在~$A$ 的某个元素，它属于 $g(x)$ 与 $\overline{A}$ 之一，而不属于另一个。正如 $\overline{Z}$ 因此不可能出现在列表 $Z_1$, $Z_2$, \dots 中一样，$\overline{A}$ 也不可能在~$g$ 的值域中。}{}

\oliflabeldef{sfr:siz:nen-alt:thm:nonenum-pownat}{将 \olref{thm:cantor} 的证明与
  \olref[nen-alt]{thm:nonenum-pownat} 的证明加以比较颇具启发意义。在那里我们证明了，对于 $\Nat$ 的任何子集列表 $N_0$, $N_1$, $N_2$, \dots，都能构造出一个集合~$D$，保证它不在该列表中。之所以能保证这一点，是因为对于每个 $n \in \Nat$，$n \in N_n$ 当且仅当 $n \notin
  D$。我们在这里遵循同样的思想，只不过现在的索引~$n$ 是~$A$ 的元素而非~$\Nat$ 的元素。集合 $D$ 的定义保证了对每个 $x \in A$，它都不同于~$g(x)$，因为 $x \in g(x)$ 当且仅当 $x \notin D$。}{}

这个证明也值得与 罗素悖论（\olref[sfr][set][rus]{thm:russells-paradox}）的证明相比较。事实上，康托尔 定理正是 罗素 本人悖论的灵感来源。
\end{explain}

\begin{prob}
  证明：对任意集合~$A$，都不可能存在单射 $g\colon \Pow{A} \to A$。提示：假设 $g\colon \Pow{A} \to A$ 是单射。考虑 $D = \Setabs{g(B)}{B \subseteq A \text{ and
  } g(B) \notin B}$。令 $x = g(D)$。利用 $g$ 为单射这一事实导出矛盾。
\end{prob}

\end{document}